% from main.tex
%%%%%%%%%%  scan idx 562  |  folio 549  %%%%%%%%%%
\textbf{Théorème 5.3.2.} L'%
application \(c_1\) (5.3.1) est un homomorphisme injectif, et l'homomorphisme
\[
 \det:\Filt^1(X)\longrightarrow\Pic(X)
\]
induit, par passage au quotient, un homomorphisme inverse à gauche de \(c_1\) :
\[
 \alpha:\Gr^1(X)\longrightarrow\Pic(X).
\]
Si toute section \(n\) du faisceau constant \(\mathbb Z_X\) est bornée (par exemple si \(X\) est «~quasi-compact~», ou si l'objet final \(e\) de \(X\) est somme finie d'objets connexes), alors \(c_1\) est un isomorphisme et \(\alpha\) son inverse.

\emph{Démonstration :} \(c_1\) est un homomorphisme, car on a pour tout \(L,L'\in\Pic(X)\) l'équation
\[
 (1-(LL')^{-1})-(1-L^{-1})-(1-{L'}^{-1})
 =
 -(1-L^{-1})(1-{L'}^{-1})\in\Filt^2(X).
\]
Admettons pour l'instant que
\[
 \det(\Filt^2(X))=\{1\}\in\Pic(X)
\]
i.e. que \(\alpha\) soit bien défini. On a alors
\[
 \alpha\circ c_1=\id_{\Pic(X)},
\]
car par 5.2.2 b) on a :
\[
 \det(1-L^{-1})=(\det(L-1))=L
\]
pour tout \(L\in\Pic(X)\). Puis on montre que
\[
 c_1\circ\alpha=\id_{\Gr^1(X)}
\]
lorsque toute section \(n\) de \(\mathbb Z_X\) est bornée. Comme \(\Filt^1(X)\) est le groupe abélien engendré par les éléments \(\cl^\circ(E)-n\), où \(E\) est un \(\mathcal O_X\)-
